%%% eskdx
\usepackage{eskdchngsheet}		% вставка листа регистрации изменений
\usepackage{eskdfreesize}		% вставка листов нестандартного формата

%%% математика
\usepackage{amsmath,amssymb}
\usepackage{gensymb}			% знак градуса командой \degree
\usepackage{icomma}				% запятая в десятичных дробях

%%% русский язык
\usepackage[T2A]{fontenc}		% поддержка русских букв
\usepackage[utf8]{inputenc}		% кодировка ввода utf8
\usepackage[russian]{babel}		% языки: русский
% выбор русского шрифта
\IfFileExists{pscyr.sty}{%
	\usepackage{pscyr}\renewcommand{\rmdefault}{ftm}}% ftm -- TimesNewRomanPSMT (Таймс)
	{\IfFileExists{literat.sty}{\usepackage{literat}}{}}
% решение проблемы копирования текста в буфер обмена
\IfFileExists{glyphtounicode.tex}{\input glyphtounicode.tex}{}
\IfFileExists{glyphtounicode-cmr.tex}{\input glyphtounicode-cmr.tex}{}	% from pdfx package
\pdfgentounicode=1
\usepackage{cmap}				% улучшенный поиск русских слов в полученном pdf-файле

%%% списки
\usepackage{enumitem}

%%% таблицы
\usepackage{multirow}			% команда для объединения строк
\usepackage{longtable,tabu}		% сложные и длинные таблицы
\usepackage{float}				% опция [H] для точного размещения плавающих объектов
\usepackage{setspace}			% отключает полуторный интервал в окружениях table, figure, предоставляет окружение spacing{} и команды \doublespacing и др.

%%% рисование
\usepackage[europeanresistors,americaninductors]{circuitikz} % рисование электрических схем (автоподключение tikz)
\usetikzlibrary{shapes,arrows.meta,chains} % библиотеки tikz

\usepackage[ 					% гипертекстовое оглавление в PDF
	bookmarks=true, colorlinks=true, unicode=true,
	urlcolor=black,linkcolor=black, anchorcolor=black,
	citecolor=black, menucolor=black, filecolor=black,
	hypertexnames=false % можно поставить ссылку из основного документа в Приложение
	]{hyperref}
\usepackage[russian]{cleveref}	% умные ссылки

%%% микротипографика
\usepackage[final]{microtype}[2016/05/14] % улучшает представление букв и слов в строках